\textbf{1. Kiến thức cần nhớ}
\begin{itemize}
\item \textbf{Đưa thừa số ra ngoài dấu căn}: Với $B \ge 0$:
$$\boxed{\sqrt{A^2 B} = |A|\sqrt{B} = \begin{cases} A\sqrt{B} & \text{khi } A \ge 0 \\ -A\sqrt{B} & \text{khi } A < 0 \end{cases}}$$
\item \textbf{Đưa thừa số vào trong dấu căn}:
\begin{itemize}
\item Với $A \ge 0, B \ge 0$: $\boxed{A\sqrt{B} = \sqrt{A^2 B}}$.
\item Với $A < 0, B \ge 0$: $\boxed{A\sqrt{B} = -\sqrt{A^2 B}}$.
\end{itemize}
\item \textbf{Khử mẫu của biểu thức lấy căn}: Với $A \cdot B \ge 0$ và $B \ne 0$:
$$\boxed{\sqrt{\dfrac{A}{B}} = \dfrac{\sqrt{AB}}{|B|}}$$
\item \textbf{Trục căn thức ở mẫu}:
\begin{itemize}
\item Mẫu là căn đơn: $\boxed{\dfrac{A}{\sqrt{B}} = \dfrac{A\sqrt{B}}{B}}\quad (B > 0)$.
\item Mẫu là tổng/hiệu có căn (nhân với lượng liên hợp):
$$\boxed{\dfrac{C}{\sqrt{A} \pm \sqrt{B}} = \dfrac{C(\sqrt{A} \mp \sqrt{B})}{A - B}}\quad (A \ge 0, B \ge 0, A \ne B).$$
\end{itemize}
\end{itemize}

\textbf{2. Phương pháp giải các dạng}

%% DANG 9D3?9-1
\textbf{Dạng 1. Đưa thừa số ra ngoài, vào trong dấu căn}
\begin{enumerate}
\item Phân tích biểu thức dưới dấu căn thành tích có chứa thừa số bình phương: $A^2 \cdot B$.
\item Đưa ra ngoài: $\sqrt{A^2 B} = |A|\sqrt{B}$. Chú ý xét dấu của $A$ theo điều kiện đề bài.
\item Khi đưa số âm vào trong dấu căn: giữ dấu trừ ở ngoài căn, chỉ đưa giá trị tuyệt đối vào trong căn: $-3\sqrt{2} = -\sqrt{3^2 \cdot 2} = -\sqrt{18}$.
\end{enumerate}

%% DANG 9D3?9-2
\textbf{Dạng 2. Khử mẫu của biểu thức lấy căn, trục căn thức ở mẫu}
\begin{enumerate}
\item Nhân cả tử và mẫu với thừa số thích hợp để mẫu trở thành bình phương hoặc hiệu hai bình phương (mất dấu căn ở mẫu).
\item Lượng liên hợp:
\begin{itemize}
\item Lượng liên hợp của $\sqrt{a} + \sqrt{b}$ là $\sqrt{a} - \sqrt{b}$.
\item Lượng liên hợp của $\sqrt{a} - b$ là $\sqrt{a} + b$.
\end{itemize}
\item Rút gọn phân số sau khi trục căn thức.
\end{enumerate}

%% DANG 9D3?9-3
\textbf{Dạng 3. Rút gọn biểu thức chứa căn thức bậc hai}
\begin{enumerate}
\item \textbf{Bước 1}: Tìm ĐKXĐ (nếu đề bài chưa cho).
\item \textbf{Bước 2}: Phân tích các mẫu thức thành nhân tử (đặt nhân tử chung, dùng hằng đẳng thức $x - 1 = (\sqrt{x}-1)(\sqrt{x}+1)$).
\item \textbf{Bước 3}: Quy đồng mẫu thức và thu gọn tử thức.
\item \textbf{Bước 4}: Chia đa thức tử cho mẫu để được phân thức tối giản.
\end{enumerate}

%% DANG 9D3?9-4
\textbf{Dạng 4. Bài toán tổng hợp: rút gọn, tính giá trị, tìm $x$}
\begin{enumerate}
\item \textbf{Tính giá trị biểu thức khi biết $x$}: Thu gọn biểu thức trước $\rightarrow$ biến đổi giá trị $x$ về dạng bình phương $(\sqrt{a} \pm \sqrt{b})^2$ để tính $\sqrt{x}$ $\rightarrow$ thay vào biểu thức đã rút gọn.
\item \textbf{Tìm $x$ để biểu thức thỏa mãn điều kiện}:
\begin{itemize}
\item Tìm $x$ để $P = k$ (giải phương trình).
\item Tìm $x$ để $P > k$ hoặc $P < k$ (chuyển vế xét dấu $P - k$).
\item Tìm $x$ nguyên để biểu thức $P$ nhận giá trị nguyên (chia tử cho mẫu rồi tìm ước của hằng số).
\end{itemize}
\end{enumerate}
